% Abstract (DECISIONS §F28, exact text; at most 1920 characters).
\begin{abstract}
Solomonoff function induction (SFI) predicts labels from inputs with a universal prior over random functions. We study the proposal to penalise program length more steeply than log-likelihood, by a factor $\lambda>1$, for hypotheses with private randomness at each input, so that learning selects a simple model plus noise rather than a model of a teacher's systematic errors. A standard objection is that a universal mixture, taken as one hypothesis, absorbs all structure. We prove that it does in every class containing its own mixture, such as Solomonoff induction and SFI with shared randomness, while no private-randomness hypothesis is universal. For a random parity on uniform inputs whose labels are flipped, at a rate below $0.056$, on an error set of product form that is incompressible relative to the parity, we prove that for $\lambda=cn$ ($n$ samples) and $c$ in an explicit window, every near-optimal hypothesis of the population objective knows the parity and gives the erroneous labels far less probability than a memoriser does. The window needs $\lambda$ proportional to $n$; density and description length of the error rule do not suffice. At $\lambda=1$ a computable sampler and the programs that fill in its sampled bits get similar posterior weight; for fixed $\lambda>1$, private randomness selects samplers over fill-ins, while with shared randomness a universal mixture beats every sufficiently long sampler. For a random parity on $\{0,1\}^d$ whose labels are flipped on a random set of inputs of density below $1/4$, the sophistication defined with product models exceeds the set sophistication of algorithmic statistics by $d-O(\log d)$ bits, for a range of sufficiency slacks.
\end{abstract}
